\documentclass[a4paper,12pt]{article}
%-------------------------------------%
%                                     %
%    parcial mhs                      %
%   fisica                 2n btx     %
%                                     %
%-------------------------------------%
\usepackage[utf8]{inputenc}
\usepackage[catalan]{babel}
\usepackage{amssymb}
\usepackage[symbol]{footmisc}
\usepackage{multicol}
\usepackage{amsmath}
\usepackage{cancel}
\usepackage{xcolor}
\usepackage{graphicx}
\usepackage[most]{tcolorbox}
\usepackage{tikz}
\usepackage{pgfplots}
\usepackage{eurosym}
\usepackage{wasysym}
\usepackage{hyperref}
\usepackage{scalerel}
\usepackage{xlop}
\usepackage{circuitikz}
\usetikzlibrary{decorations.pathmorphing}
\usetikzlibrary{arrows,positioning,shapes,fit,calc}
\usepackage{hyperref}
\hypersetup{hidelinks,}
\usepackage{fancyhdr}
\pagestyle{fancy}
\fancyhf{}
\rhead{\footnotesize A.Arroyo (aarroyo+batx@stjosep.org)}
\lhead{\footnotesize Parcial Ones 2n Batxillerat 2024-25}
\lfoot{\href{http://creativecommons.org/licenses/by-nc-nd/4.0/}{\includegraphics[width=1.5cm]{cc}}}
\rfoot{\thepage}
\setlength\headheight{14pt}
\pgfplotsset{compat=1.13}
\setlength{\arrayrulewidth}{0.3mm}
\setlength{\tabcolsep}{18pt}
\renewcommand{\arraystretch}{1.5}
%%%%%%%%%%%%%%%%%%%%%%%%%%%%%%%%%%%%%%%
\begin{document}
% \begin{enumerate}
% \item
% \begin{enumerate}
% \item
\noindent\textbf{1.} $(a)$
A partir de l'equació de l'ona
$$y(x,t)=0,3\cos\left(\frac{\pi}{4}t-\frac{\pi}{5}x+\frac{\pi}{2} \right)$$
podem trobar directament
$$A=0,3\,m\quad \omega=\frac{\pi}{4}\,rad/s\rightarrow T=\frac{2\pi}{\omega}=\frac{2\bcancel{\pi}}{\frac{\bcancel{\pi}}{4}}=8\,s$$
$$f=\frac{1}{T}=\frac{1}{8}=0,125\,s\quad \lambda=\frac{2\pi}{k}=\frac{2\bcancel{\pi}}{\frac{\bcancel{\pi}}{5}}=10\,m$$

\bigskip
\noindent$(b)$ 
Fent servir la relació $\lambda=vT$ tenim
$$v=\frac{\lambda}{T}=\frac{10}{8}=1,25\,m/s$$

\bigskip
\noindent$(c)$
La velocitat transversal dels punts de l'ona es poden calcular a partir de la derivada respecte el temps de l'equació d'ona
$$v_{y}(x,t)=\dot{y}(x,t)=-0,3\cdot\frac{\pi}{4}\sin\left(\frac{\pi}{4}t-\frac{\pi}{5}x+\frac{\pi}{2}\right)$$
llavors
\begin{equation*}
\begin{split}
v_{y}(5,20)&=-0,3\cdot\frac{\pi}{4} 
 \sin\left(\frac{\pi}{4}\cdot20-\frac{\pi}{5}\cdot 5+\frac{\pi}{2}\right)\\
 &=-0,3\cdot\frac{\pi}{4}\sin\left(4\pi+\frac{\pi}{2}\right)\\
 &=-0,3\cdot\frac{\pi}{4}\sin\left(\frac{\pi}{2}\right)=-0,3\cdot\frac{\pi}{4}=-0,2356\,m/s
\end{split}
\end{equation*}
\textbf{2.} $(a)$ A partir de la definició de intensitat sonora
$$\beta=10\log\frac{I}{I_{0}}$$
i fent servir les dades del problema podem calcular
$$180=10\log\frac{I}{I_{0}}\rightarrow I=I_{0}\cdot10^{18}=10^{-12}\cdot10^{18}=10^{6}\,W/m^{2}$$
\bigskip

\noindent$(b)$ Calculem directament a partir de la relació entre la potència, la intensitat i la distància
$$P=IA=I\cdot4\pi r^{2}=4\cdot10^{6}\cdot\pi\cdot1^{2}=1,257\cdot10^{7 }\,W$$

%\clearpage
\bigskip

\noindent$(c)$ Com que es tracta d'una ona tridimensional, la relació entre la intensitat a distàncies diferents del focus emissor es pot escriure com
$$I_{1}r^{2}_{1}=I_{2}r^{2}_{2}$$
d'on
$$r_{2}=r_{1}\sqrt{\frac{I_{1}}{I_{2}}}=1\cdot\sqrt{\frac{10^{6}}{1}}=10^{3}\,m$$
\bigskip

\noindent\textbf{3.} 
\noindent $(a)$ En l'esquema que correspon a una corda subjecta pels seus extrems vibrant en el quart harmònic es pot comprovar que en la longitud de la corda ``hi caben'' dues longituds d'ona, de forma que tenim
$$2 \lambda=1\rightarrow \lambda=0,5\,m$$
\bigskip

\noindent $(b)$
A partir de la relació $\lambda=v T$, tenim
$$v=\frac{\lambda}{T}=\lambda f=0,5\cdot 925=462,5\,m/s$$
\bigskip

\noindent $(c)$
Al vibrar en l'harmònic fonamental es comprova que en la longitud de la corda ``hi cap'' mitja longitud d'ona, llavors $\lambda=2\,m$ i la freqüència corresponent serà
$$f=\frac{v}{\lambda}=\frac{462,5}{2}=231,25\,Hz$$
on hem suposat que, per les ones estacionàries, la velocitat de fase és sempre la mateixa per tots els harmònics.
\clearpage

\noindent\textbf{4.} $(a)$
Hem de demanar que l'amplitud efectiva sigui zero, llavors
$$0,040\sin(5,0\pi x)=0\rightarrow \sin(5,0\pi x)=0$$
d'on
$$5,0\pi x=n\pi, \quad n\in \mathbb{Z}$$
ja que el sinus d'un angle val zero quan el seu argument és $0,\, \pi,\,2\pi,\, 3\pi...$. Tindrem doncs
$$5,0\bcancel{\pi} x= n\bcancel{\pi}\rightarrow x_{n}=\frac{n}{5,0}\quad n\in\mathbb{Z}$$
i finalment, 
$$x_{0}=\frac{0}{5,0}=0\,m,\quad x_{1}=\frac{1}{5,0}=0,2\,m\quad x_{2}=\frac{2}{5,0}=0,4\,m$$
i ja els tenim tots.
\bigskip

\noindent$(b)$ El període es pot trobar directament a partir de l'equació de l'ona estacionària ja que és $\omega=40\pi\,rad/s$ i tindrem
$$T=\frac{2\pi}{\omega}=\frac{2\bcancel{\pi}}{40\bcancel{\pi}}=0,05\,s$$
\bigskip

\noindent$(c)$ Com que tenim que $k=5,0\pi\,rad/m$ podem calcular
$$\lambda=\frac{2\pi}{k}=\frac{2\bcancel{\pi}}{5,0\bcancel{\pi}}=0,4\,m$$
i la velocitat demanada valdrà
$v=\frac{\lambda}{T} =\frac{0,4}{0,05}=8,0\,m/s$

\end{document}
